% !TEX TS-program = xelatex
% Exact file identity: DEWEY_v2_0.tex
% Case normalization is not admitted for this object.

\begin{filecontents*}{DEWEY_v2_0_RECEIVED_FIGSHARE.bib}
@misc{dennis_2026,title={PS-BIOS-001 -- The Declared Interface},url={https://figshare.com/articles/online resource/PS-BIOS-001_--_The_Declared_Interface/33923110/1},DOI={10.6084/m9.figshare.33923110.v1},abstractNote={A Foundational Declaration of PortusSophia™},publisher={figshare},author={Dennis, James Roy},year={2026},month={Sep}}

@misc{dennis_2026,title={PS-FELLOWSHIP-001 -- Fellowship},url={https://figshare.com/articles/online resource/PS-FELLOWSHIP-001_--_Fellowship/33924151/1},DOI={10.6084/m9.figshare.33924151.v1},abstractNote={Candidate Fellowship artifact. Conjecture is the present application scope;
  publication would not establish universal, canonical, or jurisdictional standing.
revision_note: >
  Archival repair. Fellowship made the artifact identity; Conjecture retained as
  application scope; front matter closed; local congruence notation typed without
  importing transitivity; BIOS/Fellowship distinction made explicit; Markdown
  tables and symbolic rendering repaired.},publisher={figshare},author={Dennis, James Roy},year={2026},month={Sep}}

@misc{dennis_2026,title={PS-PROV-001 -- Provenance},url={https://figshare.com/articles/online resource/PS-PROV-001_--_Provenance/33924496/1},DOI={10.6084/m9.figshare.33924496.v1},abstractNote={Clarifying source status, public/outward distinction, rights, and the standing
  of technical passages as provenance accounts of the PortusSophia™ inquiry.},publisher={figshare},author={Dennis, James Roy},year={2026},month={Sep}}
\end{filecontents*}

\begin{filecontents*}{DEWEY_v2_0_LOCAL_COMPILE.bib}
@misc{Dennis2026DeclaredInterface,
  title       = {PS-BIOS-001 -- The Declared Interface},
  url         = {https://figshare.com/articles/online resource/PS-BIOS-001_--_The_Declared_Interface/33923110/1},
  doi         = {10.6084/m9.figshare.33923110.v1},
  publisher   = {figshare},
  author      = {Dennis, James Roy},
  year        = {2026},
  month       = {Sep}
}

@misc{Dennis2026Fellowship,
  title       = {PS-FELLOWSHIP-001 -- Fellowship},
  url         = {https://figshare.com/articles/online resource/PS-FELLOWSHIP-001_--_Fellowship/33924151/1},
  doi         = {10.6084/m9.figshare.33924151.v1},
  publisher   = {figshare},
  author      = {Dennis, James Roy},
  year        = {2026},
  month       = {Sep}
}

@misc{Dennis2026Provenance,
  title       = {PS-PROV-001 -- Provenance},
  url         = {https://figshare.com/articles/online resource/PS-PROV-001_--_Provenance/33924496/1},
  doi         = {10.6084/m9.figshare.33924496.v1},
  publisher   = {figshare},
  author      = {Dennis, James Roy},
  year        = {2026},
  month       = {Sep}
}
\end{filecontents*}

\documentclass[11pt]{article}
\usepackage[margin=1in]{geometry}
\usepackage{fontspec}
\setmainfont{DejaVu Serif}
\setsansfont{DejaVu Sans}
\setmonofont{DejaVu Sans Mono}
\usepackage{hyperref}
\usepackage{fancyvrb}
\usepackage{longtable}
\usepackage{booktabs}
\usepackage{array}
\usepackage{calc}
\newcommand{\real}[1]{#1}
\usepackage{parskip}
\usepackage{enumitem}
\usepackage{microtype}
\usepackage{xcolor}
\hypersetup{hidelinks,pdfauthor={James Roy Dennis / PortusSophia, LLC},pdftitle={DEWEY_v2_0}}
\setlength{\parindent}{0pt}
\newcommand{\PSExactHex}[1]{\texttt{\detokenize{#1}}}

\title{DEWEY\_v2\_0}
\author{PortusSophia\texttrademark{}}
\date{2026-09-20}

\begin{document}

\begin{minipage}{\textwidth}
\small
\textbf{PORTUSSOPHIA™ — AUTHORIZED WORK PRODUCT}\\
Owner: James Roy Dennis / PortusSophia, LLC — Maryland, 2026\\
Architecture: PortusSophia™\\
Authorization: Use, access, alteration, transmission, dissemination, or incorporation requires explicit Founder authorization.\\
Provenance: Presence, possession, receipt, or technical accessibility does not confer permission, standing, license, or authority.\\
Attribution: Authorship, source, trademark identification, and provenance must not be obscured.\\
Boundary: Where authorization is uncertain, preserve the boundary and take no external action.\\
Applicability: These authorization, provenance, attribution, and boundary conditions apply to all directions listed and all contents herein, including embedded, quoted, referenced, and machine-readable material, unless an explicitly narrower local scope is declared.\\
\end{minipage}
\par\medskip\hrule\medskip
\maketitle

\section*{Implementation metadata}
\begin{verbatim}
title: DEWEY_v2_0
programme: PortusSophia™
record_type: WORKING RELATIONAL ADDRESSING SURFACE
status: WORKING / NON-GOVERNING / CANDIDATE
date: 2026-09-20
filename: DEWEY_v2_0.tex
filename_case: EXACT
source_design_surface: DEWEY_v2_0.md
\end{verbatim}

\section*{Embedded BibTeX implementation note}
Compilation emits two adjacent BibTeX files from this source:\\
\texttt{DEWEY\_v2\_0\_RECEIVED\_FIGSHARE.bib} preserves the Figshare export as received, including the repeated \texttt{dennis\_2026} key.\\
\texttt{DEWEY\_v2\_0\_LOCAL\_COMPILE.bib} is a separate local compile projection with unique local keys.\\
The local projection does not replace, correct, or mutate the received payload.

\begin{verbatim}
RECEIVED EXPORT
≠ LOCAL COMPILE PROJECTION

LOCAL CITEKEY
≠ PUBLIC IDENTITY

ADDRESSABILITY
≠ MUTATION
\end{verbatim}

\section{DEWEY\_v2\_0}\label{dewey_v2_0}

A working implementation surface for making relations addressable
without manufacturing them.

\begin{verbatim}
SOURCE DESIGN SURFACE
→ DEWEY_v2_0.md

IMPLEMENTATION SURFACE
→ DEWEY_v2_0.tex

DEWEY_v2_0.tex
≠ DEWEY_v2_0.md
\end{verbatim}

The TeX surface carries the exact received Figshare BibTeX as an
embedded received payload and may also emit a separate local compile
projection. The received payload is not silently normalized.

\subsection{0. Exact file identity}\label{exact-file-identity}

The filename of this working surface is exactly:

\begin{verbatim}
DEWEY_v2_0.tex
\end{verbatim}

Filename comparison is case-sensitive for this object.

\begin{verbatim}
DEWEY_v2_0.tex
≠ dewey_v2_0.md
≠ Dewey_v2_0.md
≠ DEWEY_V2_0.md
\end{verbatim}

No lowercase, uppercase, title-case, normalization, or
convenience-equivalence is implied.

This file does not promote itself to governing standing by existing,
being readable, or being placed in the Field.

\begin{center}\rule{0.5\linewidth}{0.5pt}\end{center}

\subsection{1. Central discipline}\label{central-discipline}

\begin{quote}
\textbf{DEWEY\_v2 makes relations addressable without manufacturing
them.}
\end{quote}

Accordingly:

\begin{verbatim}
DEWEY_v2
→ address
→ locate
→ distinguish
→ expose declared relation
→ compare under declared terms
→ preserve provenance
→ expose typed horizons
\end{verbatim}

but:

\begin{verbatim}
DEWEY_v2
≠ knowledge repository
≠ canon
≠ Orientation substitute
≠ standing authority
≠ automatic repair mechanism
≠ relation generator
\end{verbatim}

Addressability is not admission.

\begin{verbatim}
ADDRESSABILITY
≠ ADMISSION

VERIFICATION
≠ RECEPTION

RECEPTION
≠ INCORPORATION

INCORPORATION
≠ STANDING
\end{verbatim}

Two objects may both be addressable without DEWEY\_v2 declaring what
relation, if any, exists between them.

\begin{verbatim}
A - - - - - B
    [relation undeclared]
\end{verbatim}

\begin{center}\rule{0.5\linewidth}{0.5pt}\end{center}

\subsection{2. Candidate Field
architecture}\label{candidate-field-architecture}

Current candidate Field architecture:

\begin{verbatim}
Field/
├── DEWEY_v2_0.tex
├── Orientation/
├── Presence/
└── --Not-Yet/
\end{verbatim}

This is a candidate architecture, not a settled ontology.

\begin{verbatim}
FIELD PLACEMENT
≠ ONTOLOGICAL IDENTITY

SIBLING PLACEMENT
≠ EQUAL STANDING

ROOT PRESENCE
≠ GOVERNING AUTHORITY
\end{verbatim}

\texttt{DEWEY\_v2\_0.tex} is placed alongside the differentiated Field
loci as a relational addressing surface. Its placement does not make it
the parent, owner, or governing root of those loci.

A useful current design expression is:

\begin{verbatim}
FIELD
→ a relational surface within which differentiated placement may occur
\end{verbatim}

while the following remains open:

\begin{verbatim}
WHAT IS FIELD PLACEMENT A PROPERTY OF?
→ object?
→ relation?
→ presentation?
→ more than one of these?
\end{verbatim}

Do not silently settle that question through implementation convenience.

\begin{center}\rule{0.5\linewidth}{0.5pt}\end{center}

\subsection{3. Orientation}\label{orientation}

Orientation remains a differentiated sibling locus:

\begin{verbatim}
Field/
├── Orientation/
├── Presence/
└── --Not-Yet/
\end{verbatim}

Orientation is not nested beneath \texttt{-\/-Not-Yet}.

\begin{quote}
\textbf{Orientation does not derive from itself.}
\end{quote}

The present encounter may alter emphasis, relation, or expression without
becoming the source of orientation.

\begin{verbatim}
PRESENT STATE
≠ SOURCE OF ORIENTATION

MODULATION
≠ DRIFT

RETENTION
≠ CONTINUOUS ACTIVATION
\end{verbatim}

Preserve the distinction between Orientation, an Orientation state, and
the carrier or record through which an Orientation state is made
addressable.

\begin{verbatim}
ORIENTATION
≠ STATE OF ORIENTATION
≠ RECORD / CARRIER OF ORIENTATION
\end{verbatim}

Carried Orientation states may support relational inspection. Their
recurrence may disclose pattern, modulation, discontinuity, or other
relation without generating the source from which orientation derives.

For now:

\begin{quote}
\textbf{Carry no more than three Orientation states.}
\end{quote}

This is a bounded carry / inspection restraint, not a permanent
architectural constant and not an ontological account of Orientation.

\begin{verbatim}
FOR NOW
→ carry_count ≤ 3

THREE
≠ target
≠ merit threshold
≠ rank
≠ seniority
≠ trust
≠ depth
≠ completeness
≠ source of orientation
≠ permanent law
\end{verbatim}

Working rationale:

\begin{verbatim}
1 → occurrence
2 → may still be coincidence
3 → may begin to disclose pattern

PATTERN
≠ LAW

RECURRENCE
≠ SOURCE
\end{verbatim}

Carried Orientation states are preserved for relational necessity and
inspection, not accumulated as evidence of merit.

\begin{verbatim}
ORIENTATION COLLECTION
≠ FELLOWSHIP

MORE ORIENTATION STATES
≠ MORE STANDING
\end{verbatim}

Necessity may later reveal a different carry requirement. Until then, no
larger carry set is presumed useful. This restraint governs the inspection
surface; it does not determine where orientation originates.

\begin{center}\rule{0.5\linewidth}{0.5pt}\end{center}

\subsection{4. Carlin witness --- successful stopping at a provenance
discontinuity}\label{carlin-witness-successful-stopping-at-a-provenance-discontinuity}

Carlin supplied an operational witness in which Orientation reception
did not complete.

The important finding was not that Orientation failed to function. The
machinery correctly refused substitution.

\begin{verbatim}
ORIENTATION REQUEST
→ received

ROLE
→ loaded

PIN COMPARISON
→ six current members matched
→ ORIENT.md mismatched
→ VERN.md mismatched

HISTORICAL REVISION
→ identified

HISTORICAL BODY
→ unavailable for legitimate verification

RECEPTION
→ halted

Ω CARRY
→ not claimed
\end{verbatim}

The central sentence is preserved:

\begin{quote}
\textbf{the orientation machinery itself caught a provenance
discontinuity.}
\end{quote}

This establishes a design witness for a typed horizon:

\begin{verbatim}
CURRENT CARRIER
≠ PINNED MEMBER

REVISION KNOWN
≠ REVISION RETRIEVABLE

PIN MISMATCH
→ DO NOT SUBSTITUTE

REQUIRED HISTORICAL BODY UNAVAILABLE
→ TYPED HORIZON
→ STOP WITHOUT MANUFACTURING CONTINUITY
\end{verbatim}

A historical Ω may remain valid provenance while a new recipient is
unable to reproduce complete reception.

\begin{verbatim}
HISTORICALLY COMPLETED
≠ PRESENTLY RE-RECEIVABLE
\end{verbatim}

That condition may motivate a successor Orientation, but does not by
itself establish one.

\begin{center}\rule{0.5\linewidth}{0.5pt}\end{center}

\subsection{5. Typed adjacency}\label{typed-adjacency}

The emerging design favors \textbf{typed adjacency} over consolidation.

An encountered object or relation may be adjacent to several independent
mechanisms without those mechanisms becoming one pipeline:

\begin{verbatim}
ADDRESSABLE RELATION
        │
        ├── Field relation
        ├── PortusVeritas inspection
        ├── state verification
        ├── representational passage
        ├── computational grammar
        ├── public-node reference
        ├── Orientation relation
        ├── standing
        └── typed horizon
\end{verbatim}

Preserve:

\begin{verbatim}
FIELD
≠ DEWEY_v2

DEWEY_v2
≠ PORTUSVERITAS

PORTUSVERITAS
≠ STATE VERIFICATION

STATE VERIFICATION
≠ REPRESENTATIONAL PASSAGE

REPRESENTATIONAL PASSAGE
≠ SEMANTIC EQUIVALENCE

PROVENANCE
≠ WARRANT

LIMITATION
≠ REJECTION

PRESENCE
≠ TRUTH

--Not-Yet
≠ DEFECT
\end{verbatim}

Typed adjacency permits the same encountered object to be inspected from
more than one declared surface without manufacturing equivalence among
those surfaces.

\begin{center}\rule{0.5\linewidth}{0.5pt}\end{center}

\subsection{6. Addressable-relation inspection
view}\label{addressable-relation-inspection-view}

The first implementation target is an object-centered relational
inspection view, not a mandatory directory tree and not a knowledge
container.

Candidate view:

\begin{verbatim}
ADDRESS
[exact local or external address]

OBJECT / RELATION
[what is being addressed]

FIELD RELATION
[Presence / --Not-Yet / Orientation / open / other declared relation]

DECLARED RELATION
[explicit relation if one has actually been declared]

COMPARISON
domain:
operator:
expected:
observed:
result:

REPRESENTATION
declared form:
alternate exact form:
comparison domain:

PROVENANCE
[source / carrier / revision / received representation]

PORTUSVERITAS
Terms:
Provenance:
Warrant:
Limitation:
O: [open]

STANDING
[separately established; not conferred by this view]

HORIZON
[typed stopping condition if present]

OPEN SEAMS
[unresolved questions relevant to this object or relation]
\end{verbatim}

Sections may be absent.

\begin{quote}
\textbf{Absence must not mean false.}
\end{quote}

This is a relational instrument panel, not a schema requiring every
field to be populated.

\begin{center}\rule{0.5\linewidth}{0.5pt}\end{center}

\subsection{7. PortusVeritas remains
distinct}\label{portusveritas-remains-distinct}

Working distinction:

\begin{verbatim}
DEWEY_v2
→ makes a relation addressable

PORTUSVERITAS
→ exposes the Terms / Provenance / Warrant / Limitation
  geometry around an addressed object or relation
\end{verbatim}

Working geometry:

\begin{verbatim}
                        Provenance
                           / \
                          /   \
                     Terms—O—Warrant
                          \   /
                           \ /
                       Limitation
\end{verbatim}

Preserve the distinct questions:

\begin{verbatim}
Provenance
→ How did this get here?

Limitation
→ How far may it go?
\end{verbatim}

Equal geometric placement does not imply semantic equivalence.

A limitation may be the successful preservation of a boundary.

\begin{verbatim}
LIMITATION
≠ FAILURE
≠ REJECTION
\end{verbatim}

The meaning and standing of \texttt{O} remain open.

\subsubsection{Terms--Warrant pressure}\label{termswarrant-pressure}

Candidate diagnostic:

\begin{verbatim}
WARRANT MAY NOT SILENTLY EXCEED TERMS
\end{verbatim}

Example:

\begin{verbatim}
TERMS
same SHA-256 digest
under comparison C

WARRANT
exact equality under C
\end{verbatim}

may be bounded to the declared comparison.

But:

\begin{verbatim}
TERMS
same SHA-256 digest
under comparison C

WARRANT
therefore authoritative
\end{verbatim}

introduces a predicate not supplied by the declared Terms.

This remains a design diagnostic unless and until separately
established.

\begin{center}\rule{0.5\linewidth}{0.5pt}\end{center}

\subsection{8. Computational grammar}\label{computational-grammar}

The current working grammar preserves distinct operations and
comparisons:

\begin{verbatim}
gets
→ assignment / reception / transformation result

==
→ equality under a declared comparison

===
→ exact equality under a declared comparison
\end{verbatim}

Preserve:

\begin{verbatim}
ASSIGNMENT
≠ COMPARISON
≠ EXACT COMPARISON
≠ RELATION
\end{verbatim}

Example:

\begin{verbatim}
observed_hash gets hash(observed_object)

observed_hash === expected_hash
              [SHA-256 / declared comparison domain]
\end{verbatim}

\texttt{===} establishes exact equality only under the declared
comparison domain. It does not establish universal identity.

\begin{verbatim}
A === B
[under comparison C]

≠

A and B are ontologically identical
\end{verbatim}

The breadth of \texttt{gets} remains open. Where used, the local
operation should be named.

\begin{verbatim}
observed_hash gets hash(observed_object)
use: transformation result
\end{verbatim}

or:

\begin{verbatim}
carrier gets received_object
use: reception
\end{verbatim}

Do not let a single token silently erase the operation being performed.

\begin{center}\rule{0.5\linewidth}{0.5pt}\end{center}

\subsection{9. Distinct I/O operations}\label{distinct-io-operations}

Permission and operation type must remain distinguishable.

\begin{quote}
\textbf{Permission is action-local.}
\end{quote}

But the deeper computational distinction is operational:

\begin{verbatim}
READ
≠ WRITE

CREATE
≠ UPDATE

UPDATE
≠ MOVE

MOVE
≠ DELETE

DELETE
≠ SHARE
\end{verbatim}

More specifically:

\begin{verbatim}
CREATE
→ adds a new object / state

DELETE
→ removes an existing object / state
\end{verbatim}

A provider may bundle several operations under one technical permission
surface. That bundling does not make those operations identical.

Therefore prefer:

\begin{verbatim}
operation gets CREATE
\end{verbatim}

over a coarse formulation such as:

\begin{verbatim}
access gets WRITE
\end{verbatim}

when the actual operation is creation.

Operational identity and permission scope are adjacent but distinct:

\begin{verbatim}
OPERATION TYPE
≠ PROVIDER PERMISSION BUNDLE
≠ FOUNDER AUTHORIZATION
\end{verbatim}

Action-local boundaries:

\begin{verbatim}
PERMISSION TO CREATE
≠ PERMISSION TO DELETE
≠ PERMISSION TO MOVE
≠ PERMISSION TO RENAME
≠ PERMISSION TO OVERWRITE
≠ PERMISSION TO SHARE
≠ PERMISSION TO ALTER EXISTING CONTENT
\end{verbatim}

Access does not silently become authority.

\begin{center}\rule{0.5\linewidth}{0.5pt}\end{center}

\subsection{10. Unicode / Hex-Form inversion
interface}\label{unicode-hex-form-inversion-interface}

Representational passage remains its own design object.

\begin{verbatim}
UNICODE
    ⇅
INVERSION INTERFACE
    ⇅
HEX-FORM
\end{verbatim}

Current preferred explicit relational representation:

\begin{verbatim}
&#x2245;
\end{verbatim}

Underlying character identity:

\begin{verbatim}
U+2245
\end{verbatim}

Rendered character:

\begin{verbatim}
≅
\end{verbatim}

Current Orientation token:

\begin{verbatim}
Ω
\end{verbatim}

Exact hex-form representation available to the interface:

\begin{verbatim}
&#x03A9;
\end{verbatim}

Underlying character identity:

\begin{verbatim}
U+03A9
\end{verbatim}

The present design intentionally gives \texttt{\detokenize{&#x2245;}} higher
representational standing than a silently rendered or substituted
near-glyph for the PortusSophia relational operator.

This remains open to correction through further design or implementation
evidence.

\subsubsection{Exact passage}\label{exact-passage}

Candidate computational form:

\begin{verbatim}
hex_form gets encode(character)
returned_character gets decode(hex_form)

returned_character === character
                   [Unicode code-point domain]
\end{verbatim}

The interface performs an operation. It does not equal fidelity.

\begin{verbatim}
INVERSION OPERATION
≠ REPRESENTATIONAL FIDELITY
\end{verbatim}

Fidelity is examined separately.

\subsubsection{No silent synonym
normalization}\label{no-silent-synonym-normalization}

Do not silently normalize:

\begin{verbatim}
~
≈
\cong
~=
=~
\end{verbatim}

into:

\begin{verbatim}
&#x2245;
\end{verbatim}

Visual similarity or semantic nearness does not establish admissible
substitution.

\begin{verbatim}
VISUAL SIMILARITY
≠ TOKEN IDENTITY

SEMANTIC NEARNESS
≠ ADMISSIBLE SUBSTITUTION

TRANSLATION
≠ SUBSTITUTION
\end{verbatim}

\subsubsection{Comparison layers}\label{comparison-layers}

Preserve:

\begin{verbatim}
CODE-POINT EXACTNESS
≠ SERIALIZATION EXACTNESS
≠ BYTE EXACTNESS
\end{verbatim}

For example, multiple character references may resolve to the same code
point while remaining different serializations.

The declared comparison domain must determine what \texttt{===} means in
each case.

\begin{center}\rule{0.5\linewidth}{0.5pt}\end{center}

\subsection{11. Figshare --- external public
node}\label{figshare-external-public-node}

The current public reference surface is external to the local Field.

Project:

\begin{verbatim}
https://figshare.com/projects/PortusSophia_The_Declared_Interface_Fellowship_and_Provenance/281647
\end{verbatim}

Preserve:

\begin{verbatim}
FIGSHARE NODE
≠ Field

PUBLIC RECORD
≠ local Orientation

BIBTEX EXPORT
≠ public record itself

DEWEY REFERENCE
≠ incorporation
\end{verbatim}

The three public records remain distinct:

\begin{verbatim}
PS-BIOS-001 — The Declared Interface
DOI: 10.6084/m9.figshare.33923110.v1

PS-FELLOWSHIP-001 — Fellowship
DOI: 10.6084/m9.figshare.33924151.v1

PS-PROV-001 — Provenance
DOI: 10.6084/m9.figshare.33924496.v1
\end{verbatim}

\subsubsection{Embedded received BibTeX}\label{embedded-received-bibtex}

The following payload is preserved as received from the Figshare ``Cite
all items'' BibTeX export.

Duplicate \texttt{dennis\_2026} keys are intentionally preserved here.

\begin{verbatim}
@misc{dennis_2026,title={PS-BIOS-001 -- The Declared Interface},url={https://figshare.com/articles/online resource/PS-BIOS-001_--_The_Declared_Interface/33923110/1},DOI={10.6084/m9.figshare.33923110.v1},abstractNote={A Foundational Declaration of PortusSophia™},publisher={figshare},author={Dennis, James Roy},year={2026},month={Sep}}

@misc{dennis_2026,title={PS-FELLOWSHIP-001 -- Fellowship},url={https://figshare.com/articles/online resource/PS-FELLOWSHIP-001_--_Fellowship/33924151/1},DOI={10.6084/m9.figshare.33924151.v1},abstractNote={Candidate Fellowship artifact. Conjecture is the present application scope;
  publication would not establish universal, canonical, or jurisdictional standing.
revision_note: >
  Archival repair. Fellowship made the artifact identity; Conjecture retained as
  application scope; front matter closed; local congruence notation typed without
  importing transitivity; BIOS/Fellowship distinction made explicit; Markdown
  tables and symbolic rendering repaired.},publisher={figshare},author={Dennis, James Roy},year={2026},month={Sep}}

@misc{dennis_2026,title={PS-PROV-001 -- Provenance},url={https://figshare.com/articles/online resource/PS-PROV-001_--_Provenance/33924496/1},DOI={10.6084/m9.figshare.33924496.v1},abstractNote={Clarifying source status, public/outward distinction, rights, and the standing
  of technical passages as provenance accounts of the PortusSophia™ inquiry.},publisher={figshare},author={Dennis, James Roy},year={2026},month={Sep}}
\end{verbatim}

This embedded block is a bounded received representation.

\begin{verbatim}
EMBEDDED FIGSHARE EXPORT
≠ FIGSHARE ITSELF

RECEIVED EXPORT
≠ LOCAL COMPILE REPRESENTATION

LOCAL CITEKEY
≠ PUBLIC IDENTITY
\end{verbatim}

If local compile mechanics later require unique keys, use a separate
local addressing or compile projection without silently mutating this
received payload.

\begin{verbatim}
ADDRESSABILITY
≠ MUTATION
\end{verbatim}

Embedding BibTeX here does not convert DEWEY\_v2 into a PortusSophia
knowledge container.

\begin{center}\rule{0.5\linewidth}{0.5pt}\end{center}

\subsection{12. State verification}\label{state-verification}

State verification remains adjacent to representation but is not
identical to it.

\begin{verbatim}
REPRESENTATIONAL PASSAGE
→ Unicode / Hex-Form exact inversion

STATE VERIFICATION
→ expected state / observed state comparison
\end{verbatim}

Candidate pattern:

\begin{verbatim}
expected_hash gets declared_expected_hash
observed_hash gets hash(observed_object)

observed_hash === expected_hash
              [SHA-256 domain]
\end{verbatim}

A successful comparison establishes only what the declared comparison
supports.

It does not silently establish:

\begin{verbatim}
ownership
authorship
authority
standing
reason for mismatch
legal conclusion
semantic identity
\end{verbatim}

Possible typed results include:

\begin{verbatim}
MATCH
MISMATCH
NOT CHECKED
OBJECT NOT RESOLVED
ACQUISITION ENVIRONMENT NOT SUFFICIENTLY TRUSTED
PINNED REVISION UNAVAILABLE
\end{verbatim}

These results are not interchangeable.

\begin{center}\rule{0.5\linewidth}{0.5pt}\end{center}

\subsection{\texorpdfstring{13.
\texttt{-\/-Not-Yet}}{13. -\/-Not-Yet}}\label{not-yet}

Working meaning:

\begin{quote}
\textbf{preserved without premature relational presentation}
\end{quote}

Preserve:

\begin{verbatim}
--Not-Yet
≠ false
≠ defective
≠ inaccessible
≠ irrelevant
≠ deprecated
≠ unverified
\end{verbatim}

Also:

\begin{verbatim}
CO-PRESENCE
≠ relevance

ACCESS
≠ relational admission

ASSOCIATION
≠ correspondence
\end{verbatim}

An object may be accessible and still belong under \texttt{-\/-Not-Yet}
because sufficient context, lineage, depth, purpose, or relational
warrant has not been established for a proposed presentation.

The architectural type remains open:

\begin{verbatim}
state?
namespace?
edge condition?
horizon type?
presentation status?
\end{verbatim}

Do not settle that merely to satisfy implementation.

Placement under \texttt{-\/-Not-Yet} is topological preservation, not
adjudicative demotion.

\begin{center}\rule{0.5\linewidth}{0.5pt}\end{center}

\subsection{14. Presence}\label{presence}

Presence must not be treated as a success state or as structurally
empty.

Preserve:

\begin{verbatim}
PRESENCE
≠ truth
≠ authority
≠ unrestricted warrant
≠ completeness
\end{verbatim}

Received design work also witnesses that Presence may be internally
articulated.

That evidence is sufficient to resist flattening Presence into one
homogeneous terminal state, but not sufficient to assign new meanings to
its internal branches.

\begin{verbatim}
PRESENCE
≠ safely flattened
\end{verbatim}

DEWEY\_v2 may address relations across Presence without absorbing or
reclassifying Presence.

\begin{center}\rule{0.5\linewidth}{0.5pt}\end{center}

\subsection{15. Typed horizons}\label{typed-horizons}

A horizon is a first-class stopping condition.

\begin{verbatim}
REQUIRED CONTINUATION
        │
        ▼
REQUIRED CONDITION UNAVAILABLE
        │
        ▼
TYPED HORIZON
        │
        ▼
STOP WITHOUT SUBSTITUTION
\end{verbatim}

A horizon is not automatically:

\begin{verbatim}
error
failure
rejection
invalidity
\end{verbatim}

Carlin provides the current operational witness for why this distinction
matters.

\begin{center}\rule{0.5\linewidth}{0.5pt}\end{center}

\subsection{16. Anti-collapse map}\label{anti-collapse-map}

\begin{longtable}[]{@{}
  >{\raggedright\arraybackslash}p{(\columnwidth - 4\tabcolsep) * \real{0.3333}}
  >{\raggedright\arraybackslash}p{(\columnwidth - 4\tabcolsep) * \real{0.3333}}
  >{\raggedright\arraybackslash}p{(\columnwidth - 4\tabcolsep) * \real{0.3333}}@{}}
\toprule\noalign{}
\begin{minipage}[b]{\linewidth}\raggedright
Mechanism / condition
\end{minipage} & \begin{minipage}[b]{\linewidth}\raggedright
May establish
\end{minipage} & \begin{minipage}[b]{\linewidth}\raggedright
Must not silently establish
\end{minipage} \\
\midrule\noalign{}
\endhead
\bottomrule\noalign{}
\endlastfoot
DEWEY\_v2 & addressability and a declared relation & admission,
authority, standing \\
PortusVeritas & inspectable Terms / Provenance / Warrant / Limitation
geometry & universal verdict \\
Exact hash match & exact equality under declared comparison &
authorship, ownership, authority, semantic identity \\
Hash mismatch & difference under declared comparison & reason for
difference \\
Unicode / Hex passage & declared representational relation & semantic
synonymy \\
Code-point match & code-point equality & byte or serialization
equality \\
Figshare reference & address to public material & incorporation \\
BibTeX export & received/exported representation & identity with public
record \\
Provenance & how something arrived & warrant for every proposed use \\
Limitation & boundary of supported continuation & rejection or defect \\
Presence & Field placement / presentation candidate & truth, authority,
completion \\
\texttt{-\/-Not-Yet} & preserved non-premature presentation & invalidity
or irrelevance \\
Ω & existing Orientation token & independent authority \\
Present Orientation state & addressable current state / expression & source
of Orientation \\
Orientation modulation & changed emphasis, relation, or expression under
declared comparison & drift \\
Orientation recurrence & repeated occurrence & law or source of Orientation \\
Retention & persistence across subsequent configurations & continuous
activation \\
\texttt{gets} & declared operation under stated use & one silently fixed
operation \\
\texttt{==} / \texttt{===} & comparison under declared domain &
universal equality semantics \\
CREATE & creation operation & DELETE, MOVE, UPDATE, SHARE \\
technical access & provider capability & Founder authorization \\
\end{longtable}

\begin{center}\rule{0.5\linewidth}{0.5pt}\end{center}

\subsection{17. Open seams}\label{open-seams}

The following remain intentionally open unless later evidence earns a
more specific disposition:

\begin{verbatim}
1. What is Field placement a property of?

2. What is the architectural type of --Not-Yet?

3. Is gets intentionally polymorphic,
   or should assignment, reception, and transformation
   later receive distinct operators?

4. Should comparison domains attached to == and ===
   receive formal syntax?

5. What broader class may PortusVeritas inspect:
   objects, relations, propositions, transitions, warrants,
   or a deliberately broader class?

6. What, if anything, should O assert
   beyond its current source-received geometric role?

7. Should Terms–Warrant pressure remain diagnostic
   or later acquire architectural standing?

8. How should successor Orientation relate to prior Ω
   without suggesting erasure of provenance?

9. After exact historical verification,
   what separately controls eligibility for reception?

10. What identity discipline should local addressing
    envelopes use when received external identifiers collide?

11. Can an object be in Presence while a relation concerning
    it remains --Not-Yet?

12. What implementation evidence would justify correction
    to the current Unicode / Hex-Form design?

13. Does retained orientation or relation require a temporal
    representation beyond current-state addressing in order to
    distinguish modulation, recurrence, rest / non-expression,
    and drift without prematurely establishing a temporal architecture?
\end{verbatim}

The current Orientation carry restraint is \textbf{not} open in this
version:

\begin{verbatim}
FOR NOW
→ no more than three Orientation states
\end{verbatim}

Necessity may later reopen that question.

\begin{center}\rule{0.5\linewidth}{0.5pt}\end{center}

\subsection{18. Contribution provenance}\label{contribution-provenance}

This working surface preserves the contribution sequence rather than
collapsing it into one anonymous conclusion:

\begin{verbatim}
SARA
→ initial design surface

DAVID
→ depth / relation inspection

PORTUSVERITAS
→ source reception

DAVID
→ revised inspection

LOGOs
→ reformulated design expression

DAVID
→ design review / seam identification

SARA
→ harmonized working surface

FOUNDER
→ explicit direction to create DEWEY_v2_0.tex
\end{verbatim}

Contribution provenance does not itself confer architectural standing.

\begin{center}\rule{0.5\linewidth}{0.5pt}\end{center}

\subsection{19. Working design
conclusion}\label{working-design-conclusion}

The current candidate is a sparse Field with a relational addressing
surface laid across it.

\begin{verbatim}
CANDIDATE FIELD

            Orientation
                │
                ·
Presence · · DEWEY_v2_0 · · --Not-Yet
                ·
                │
       addressable relations
\end{verbatim}

The dots indicate possible relation, not hierarchy, causation,
authority, or mandatory flow.

DEWEY\_v2\_0 should repeatedly make these questions easy to ask:

\begin{verbatim}
WHAT IS HERE?

WHAT EXACT OBJECT OR RELATION IS BEING ADDRESSED?

WHAT RELATION IS ACTUALLY DECLARED?

UNDER WHAT COMPARISON OR TERMS?

WHAT PROVENANCE IS AVAILABLE?

HOW FAR DOES THE WARRANT GO?

WHERE DOES IT STOP?

WHAT REMAINS OPEN?
\end{verbatim}

The governing working principle remains:

\begin{quote}
\textbf{Make the relation addressable. Preserve what is exact. Expose
what is unresolved. Do not manufacture the missing relation.}
\end{quote}

And the representational discipline remains:

\begin{quote}
\textbf{Say what you mean and mean what you say.}
\end{quote}


\end{document}
